% Generated mechanically from the frozen spaces-resolve.tex target.
% Do not edit this generated file; edit the bound locale source instead.

% OK, start here.
%

\setcounter{chapter}{88}
\chapter{再论曲面解消}



\phantomsection
\label{spaces-resolve-section-phantom}


\section{引言}
\label{spaces-resolve-section-introduction}

\noindent
本章讨论维数为 $2$ 的诺特代数空间的奇点解消。
我们在前面的一章中已经依照 Lipman \cite{Lipman}
讨论过概形的曲面解消；见《曲面解消》第
\ref{resolve-section-introduction} 节。
本章大多数结果都是概形情形相应结果的直接推论。

\medskip\noindent
除非特别说明，本章中的所有几何对象都是代数空间。因此，若我们说
``设 $f : X \to Y$ 是一个修改''，则意指 $f$ 是《域上的空间》定义
\ref{spaces-over-fields-definition-modification} 中的态射。
固有态射等术语也作同样理解。









\section{修改}
\label{spaces-resolve-section-modifications}

\noindent
设 $(A, \mathfrak m, \kappa)$ 为诺特局部环。令
$S = \Spec(A)$ 且 $U = S \setminus \{\mathfrak m\}$。本节考虑范畴
\begin{equation}
\label{spaces-resolve-equation-modification}
\left\{
f : X \longrightarrow S
\quad \middle| \quad
\begin{matrix}
X\text{ 是代数空间}\\
f\text{ 是固有态射}\\
f^{-1}(U) \to U\text{ 是同构}
\end{matrix}
\right\}.
\end{equation}
从 $X/S$ 到 $X'/S$ 的态射，是与 $S$ 上结构态射相容的代数空间态射
$X \to X'$。我们在《形式空间的代数化》第
\ref{restricted-section-modifications} 节中已经看到，此范畴只依赖于
$A$ 的完备化，并证明了其中对象的一些初等性质。本节专门研究
$\dim(A) \leq 2$，或闭纤维维数至多为 $1$ 的情形。

\begin{lemma}
\label{spaces-resolve-lemma-modification}
设 $(A, \mathfrak m, \kappa)$ 为维数 $2$ 的诺特局部整环，并且
$U = \Spec(A) \setminus \{\mathfrak m\}$ 是正规概形。
则任意修改 $f : X \to \Spec(A)$ 都是
(\ref{spaces-resolve-equation-modification}) 中的态射。
\end{lemma}

\begin{proof}
设 $f : X \to S$ 为修改。我们须证明 $f^{-1}(U) \to U$ 是同构。
由于 $U$ 的每个闭点 $u$ 都具有余维 $1$，结论由《域上的空间》引理
\ref{spaces-over-fields-lemma-modification-normal-iso-over-codimension-1}
得到。
\end{proof}

\begin{lemma}
\label{spaces-resolve-lemma-closed-immersion-on-fibre}
设 $(A, \mathfrak m, \kappa)$ 为诺特局部环。
设 $g : X \to Y$ 为范畴 (\ref{spaces-resolve-equation-modification}) 中的态射。
若特殊纤维上的诱导态射 $X_\kappa \to Y_\kappa$ 是闭浸入，
则 $g$ 是闭浸入。
\end{lemma}

\begin{proof}
这是《空间态射进阶》引理
\ref{spaces-more-morphisms-lemma-where-closed-immersion}
的一个特殊情形。
\end{proof}

\begin{lemma}
\label{spaces-resolve-lemma-dimension-special-fibre}
设 $(A, \mathfrak m, \kappa)$ 为维数 $\geq 1$ 的诺特局部整环。
设 $f : X \to \Spec(A)$ 为代数空间态射。假设下列条件中至少一个成立：
\begin{enumerate}
\item $f$ 是修改（《域上的空间》定义
\ref{spaces-over-fields-definition-modification}）；
\item $f$ 是变换（《域上的空间》定义
\ref{spaces-over-fields-definition-alteration}）；
\item $f$ 局部有限型且拟分离，$X$ 整，并且 $|X|$ 中恰有一个点
映到 $\Spec(A)$ 的泛点；
\item $f$ 局部有限型，$X$ 合宜，并且 $|X|$ 中映到
$\Spec(A)$ 泛点的点，正是 $|X|$ 各不可约分支的泛点；
% P10-E0403
\item 此处还可增加其他条件。
\end{enumerate}
则 $\dim(X_\kappa) \leq \dim(A) - 1$。
\end{lemma}

\begin{proof}
情形 (1)、(2)、(3) 都是 (4) 的特殊情形。取仿射概形
$U = \Spec(B)$ 以及平展态射 $U \to X$。环映射 $A \to B$ 有限型。
我们须证明 $\dim(U_\kappa) \leq \dim(A) - 1$。
由于 $X$ 合宜，$U$ 各不可约分支的泛点就是位于 $|X|$
各不可约分支泛点之上的点；见《合宜空间》引理
\ref{decent-spaces-lemma-decent-generic-points}。
因此，$\Spec(B) \to \Spec(A)$ 在 $(0)$ 上的纤维是 $B$ 的极小素理想
$\mathfrak q_1, \ldots, \mathfrak q_r$ 所成的有限集合。
于是对某个非零 $f \in A$，$A_f \to B_f$ 有限
（《代数》引理 \ref{algebra-lemma-generically-finite}）。
由此 $\kappa(\mathfrak q_i)$ 是 $A$ 的分式域的有限扩张。
设 $\mathfrak q \subset B$ 是位于 $\mathfrak m$ 之上的素理想。则
$$
\dim(B_\mathfrak q) = \max \dim((B/\mathfrak q_i)_{\mathfrak q})
\leq \dim(A)
$$
其中不等式来自 $A \subset B/\mathfrak q_i$ 的维数公式；见《代数》引理
\ref{algebra-lemma-dimension-formula}。
然而，$B_\mathfrak q/\mathfrak m B_\mathfrak q$ 的维数
（这正是 $U_\kappa$ 在相应点处的局部环的维数）至少少一，
因为极小素理想 $\mathfrak q_i$ 不属于 $V(\mathfrak m)$。
最后应用《性质》引理 \ref{properties-lemma-dimension}。
\end{proof}

\begin{lemma}
\label{spaces-resolve-lemma-modification-of-dim-2-is-projective-over-complete}
若 $(A, \mathfrak m, \kappa)$ 为维数 $2$ 的完备诺特局部整环，
则 $\Spec(A)$ 的每个修改都在 $A$ 上射影。
\end{lemma}

\begin{proof}
由《空间态射进阶》引理
\ref{spaces-more-morphisms-lemma-projective-over-complete}，
只须证明 $\Spec(A)$ 的任意修改 $X$ 的特殊纤维维数 $\leq 1$。
这由引理 \ref{spaces-resolve-lemma-dimension-special-fibre} 得出。
\end{proof}






\section{策略}
\label{spaces-resolve-section-strategy}

\noindent
设 $S$ 为概形，$X$ 为 $S$ 上的合宜代数空间。
设 $x_1, \ldots, x_n \in |X|$ 为两两不同的闭点。
对每个 $i$，依照《合宜空间》引理
\ref{decent-spaces-lemma-decent-space-elementary-etale-neighbourhood}，
选取初等平展邻域 $(U_i, u_i) \to (X, x_i)$。
这意味着 $U_i$ 是仿射概形，$U_i \to X$ 平展，$u_i$ 是 $U_i$
中位于 $x_i$ 之上的唯一点，并且
$\Spec(\kappa(u_i)) \to X$ 是表示 $x_i$ 的单态射。
缩小 $U_i$ 后，我们可以并且确实假设：若 $j \not = i$，
则 $U_i$ 中不存在映到 $x_j$ 的点。注意 $u_i \in U_i$ 是闭点。

\medskip\noindent
以 $\mathcal{C}_{X, \{x_1, \ldots, x_n\}}$ 表示如下代数空间态射
$f : Y \to X$ 所成的范畴：它们诱导同构
$f^{-1}(X \setminus \{x_1, \ldots, x_n\}) \to
X \setminus \{x_1, \ldots, x_n\}$。
对每个 $i$，以 $\mathcal{C}_{U_i, u_i}$ 表示如下代数空间态射
$g_i : Y_i \to U_i$ 所成的范畴：它们诱导同构
$g_i^{-1}(U_i \setminus \{u_i\}) \to U_i \setminus \{u_i\}$。
基变换定义了一个函子
\begin{equation}
\label{spaces-resolve-equation-equivalence}
F :
\mathcal{C}_{X, \{x_1, \ldots, x_n\}}
\longrightarrow
\mathcal{C}_{U_1, u_1} \times \ldots \times \mathcal{C}_{U_n, u_n}.
\end{equation}
为了把本章至少一部分问题约化到概形情形，我们有下述引理。

\begin{lemma}
\label{spaces-resolve-lemma-equivalence}
函子 $F$（\ref{spaces-resolve-equation-equivalence}）是范畴等价。
\end{lemma}

\begin{proof}
当 $n = 1$ 时，这就是《空间的极限》引理
\ref{spaces-limits-lemma-excision-modifications}。
当 $n > 1$ 时，可以用完全相同的方法证明，也可以从前一情形推出。
例如，假设 $g_i : Y_i \to U_i$ 是
$\mathcal{C}_{U_i, u_i}$ 中的对象。由 $n = 1$ 的情形，
可以找到 $f'_i : Y'_i \to X$，它在 $X \setminus \{x_i\}$ 上为同构，
而其到 $U_i$ 的基变换为 $f_i$。% P10-E0404
于是可以令
$$
f : Y = Y'_1 \times_X \ldots \times_X Y'_n \to X
$$
这是 $\mathcal{C}_{X, \{x_1, \ldots, x_n\}}$ 中的对象，
沿 $U_i \to X$ 作基变换后恢复 $g_i$。因此该函子本质满。
完全忠实性的证明从略。% P10-E0405
\end{proof}

\begin{lemma}
\label{spaces-resolve-lemma-equivalence-properties}
设 $X, x_i, U_i \to X, u_i$ 如 (\ref{spaces-resolve-equation-equivalence}) 中所述。
若在 $F$ 下 $f : Y \to X$ 对应于 $g_i : Y_i \to U_i$，则
$f$ 拟紧、拟分离、分离、局部有限呈示、有限呈示、局部有限型、有限型、
固有、整或有限，当且仅当对 $i = 1, \ldots, n$，$g_i$ 具有相应性质。
\end{lemma}

\begin{proof}
这由《空间的极限》引理
\ref{spaces-limits-lemma-excision-modifications-properties} 得出。
\end{proof}

\begin{lemma}
\label{spaces-resolve-lemma-equivalence-fibre}
设 $X, x_i, U_i \to X, u_i$ 如 (\ref{spaces-resolve-equation-equivalence}) 中所述。
若在 $F$ 下 $f : Y \to X$ 对应于 $g_i : Y_i \to U_i$，
则代数空间 $Y_{x_i} \cong (Y_i)_{u_i}$。
\end{lemma}

\begin{proof}
这是显然的，因为 $u_i \to x_i$ 是同构。
\end{proof}







\section{由二次变换支配}
\label{spaces-resolve-section-quadratic-spaces}

\noindent
只有当 $X$ 合宜时，我们才定义空间在一点处的爆破。

\begin{definition}
\label{spaces-resolve-definition-blowup-at-point}
设 $S$ 为概形，$X$ 为 $S$ 上的合宜代数空间，
$x \in |X|$ 为闭点。由《合宜空间》引理
\ref{decent-spaces-lemma-decent-space-closed-point}，
可用闭浸入 $i : \Spec(k) \to X$ 表示 $x$。
所谓 $X$ 在 $x$ 处的{\it 爆破 $X' \to X$}，
是指 $X$ 沿闭子空间 $Z = i(\Spec(k)) \subset X$ 的爆破。
\end{definition}

\noindent
在这种一般性下，$X$ 在 $x$ 处的爆破未必固有。
然而，若 $X$ 局部诺特，则由《空间上的除子》引理
\ref{spaces-divisors-lemma-blowing-up-projective}，该爆破是固有的。
回忆，局部诺特代数空间为诺特的，当且仅当它拟紧且拟分离。
此外，对局部诺特代数空间而言，拟分离等价于合宜
（《合宜空间》引理
\ref{decent-spaces-lemma-locally-Noetherian-decent-quasi-separated}）。

\begin{lemma}
\label{spaces-resolve-lemma-equivalence-sequence-blowups}
设 $X, x_i, U_i \to X, u_i$ 如 (\ref{spaces-resolve-equation-equivalence}) 中所述，
并假设在 $F$ 下 $f : Y \to X$ 对应于 $g_i : Y_i \to U_i$。
则 $f$ 存在分解
$$
Y = Z_m \to Z_{m - 1} \to \ldots \to Z_1 \to Z_0 = X
$$
其中 $Z_{j + 1} \to Z_j$ 是 $Z_j$ 在位于
$\{x_1, \ldots, x_n\}$ 之上的闭点 $z_j$ 处的爆破，
当且仅当对每个 $i$，$g_i$ 存在分解
$$
Y_i = Z_{i, m_i} \to Z_{i, m_i - 1} \to \ldots \to
Z_{i, 1} \to Z_{i, 0} = U_i
$$
其中 $Z_{i, j + 1} \to Z_{i, j}$ 是 $Z_{i, j}$
在位于 $u_i$ 之上的闭点 $z_{i, j}$ 处的爆破。
\end{lemma}

\begin{proof}
爆破是可表态射。因此在任一情形中，归纳可见
$Z_j \to X$ 或 $Z_{i, j} \to U_i$ 可表。
于是由《合宜空间》引理
\ref{decent-spaces-lemma-representable-named-properties}，
每个 $Z_j$ 或 $Z_{i,j}$ 都是合宜代数空间。
这说明上述断言有意义（因为爆破只对合宜空间定义）。

为证明等价性，先从一列爆破
$Z_m \to Z_{m - 1} \to \ldots \to Z_1 \to Z_0 = X$ 出发。
第一个态射 $Z_1 \to X$ 是对某个 $x_i$ 的爆破，不妨设为 $x_1$。
把 $F$ 应用于 $Z_1 \to X$，便得到在 $u_1$ 处的爆破
$Z_{1, 1} \to U_1$ 是在 $u_1$ 处的爆破；% P10-E0406
而对 $i > 1$，$Z_{i,0}=U_i$。
下一步，我们或者在 $Z_1$ 上爆破某个 $x_i$（$i\geq 2$），
或者选取 $Z_1 \to X$ 在 $x_1$ 上纤维中的闭点 $z_1$。
第一种情形中做法显然；第二种情形中，利用
$(Z_1)_{x_1} \cong (Z_{1, 1})_{u_1}$
（引理 \ref{spaces-resolve-lemma-equivalence-fibre}），
得到与 $z_1$ 对应的闭点 $z_{1,1}\in Z_{1,1}$。
然后令 $Z_{1,2}\to Z_{1,1}$ 为在 $z_{1,1}$ 处的爆破。
依此继续，便构造出每个 $g_i$ 的分解。

\medskip\noindent
反过来，给定各列爆破
$Z_{i, m_i} \to Z_{i, m_i - 1} \to \ldots \to Z_{i, 1}
\to Z_{i, 0} = U_i$，以完全相同的方式构造 $X$ 的爆破序列。
\end{proof}

\begin{lemma}
\label{spaces-resolve-lemma-make-ideal-principal}
设 $S$ 为概形，$X$ 为 $S$ 上的诺特代数空间。
设 $T \subset |X|$ 为由闭点 $x$ 组成的有限集合，并且
(1) $X$ 在 $x$ 处正则；(2) $X$ 在 $x$ 处的局部环维数为 $2$。
设 $\mathcal{I} \subset \mathcal{O}_X$ 为拟凝聚理想层，
使 $\mathcal{O}_X/\mathcal{I}$ 支撑于 $T$。则存在序列
$$
X_m \to X_{m - 1} \to \ldots \to X_1 \to X_0 = X
$$
其中 $X_{j + 1} \to X_j$ 是 $X_j$ 在某个闭点 $x_j$ 处的爆破，
该点位于 $T$ 中一点之上，并且
$\mathcal{I}\mathcal{O}_{X_n}$ 是可逆理想层。% P10-E0407
\end{lemma}

\begin{proof}
设 $T = \{x_1, \ldots, x_r\}$。如第 \ref{spaces-resolve-section-strategy} 节，
选取初等平展邻域 $(U_i,u_i)\to(X,x_i)$。
对每个 $i$，限制
$\mathcal{I}_i=\mathcal{I}|_{U_i}\subset\mathcal{O}_{U_i}$
是支撑于 $u_i$ 的拟凝聚理想层。
$U_i$ 在 $u_i$ 处的局部环正则且维数为 $2$。
因此可以应用《曲面解消》引理
\ref{resolve-lemma-make-ideal-principal}，得到序列
$$
X_{i, m_i} \to X_{i, m_i - 1} \to \ldots \to X_1
\to X_{i, 0} = U_i
$$
% P10-E0408
其中各态射是在位于 $u_i$ 之上的闭点处作爆破，且
$\mathcal{I}_i \mathcal{O}_{X_{i, m_i}}$ 可逆。
由引理 \ref{spaces-resolve-lemma-equivalence-sequence-blowups}，得到一列爆破
$$
X_m \to X_{m - 1} \to \ldots \to X_1 \to X_0 = X
$$
如引理所述，它到各 $U_i$ 的基变换给出上述序列。
于是 $\mathcal{I}\mathcal{O}_{X_n}$ 是可逆理想层。
事实上，在 $X\setminus\{x_1,\ldots,x_n\}$ 上这是已知的，% P10-E0409
而在每个 $x_i$ 的纤维的一个平展邻域中，这由构造成立。
\end{proof}

\begin{lemma}
\label{spaces-resolve-lemma-dominate-by-blowing-up-in-points}
设 $S$ 为概形，$X$ 为 $S$ 上的诺特代数空间。
设 $T\subset|X|$ 为由闭点 $x$ 组成的有限集合，并且
(1) $X$ 在 $x$ 处正则；(2) $X$ 在 $x$ 处的局部环维数为 $2$。
设 $f:Y\to X$ 为代数空间的固有态射，并且在
$U=X\setminus T$ 上为同构。则存在序列
$$
X_n \to X_{n - 1} \to \ldots \to X_1 \to X_0 = X
$$
其中 $X_{i+1}\to X_i$ 是 $X_i$ 在闭点 $x_i$ 处的爆破，
该点位于 $T$ 中一点之上，并且复合态射存在分解
$X_n\to Y\to X$。
\end{lemma}

\begin{proof}
由《空间态射进阶》引理
\ref{spaces-more-morphisms-lemma-dominate-modification-by-blowup}，
存在支配 $Y\to X$ 的 $U$-容许爆破 $X'\to X$。
因此可以假设存在理想层 $\mathcal{I}\subset\mathcal{O}_X$，
使 $\mathcal{O}_X/\mathcal{I}$ 支撑于 $T$，并且 $Y$
是 $X$ 沿 $\mathcal{I}$ 的爆破。
由引理 \ref{spaces-resolve-lemma-make-ideal-principal}，存在序列
$$
X_n \to X_{n - 1} \to \ldots \to X_1 \to X_0 = X
$$
其中 $X_{i+1}\to X_i$ 是 $X_i$ 在位于 $T$ 中一点之上的闭点
$x_i$ 处的爆破，并且 $\mathcal{I}\mathcal{O}_{X_n}$ 是可逆理想层。
由爆破的泛性质（《空间上的除子》引理
\ref{spaces-divisors-lemma-universal-property-blowing-up}），
得到所需分解。
\end{proof}





\section{由正规化爆破支配}
\label{spaces-resolve-section-normalized-blowups}

\noindent
本节证明，曲面的一个修改可由一列在点处的正规化爆破支配。

\begin{definition}
\label{spaces-resolve-definition-normalized-blowup}
设 $S$ 为概形，$X$ 为 $S$ 上的合宜代数空间，并满足
《空间的态射》引理
\ref{spaces-morphisms-lemma-prepare-normalization} 中的等价条件。
设 $x\in|X|$ 为闭点。所谓 $X$ 在 $x$ 处的
{\it 正规化爆破}，是复合 $X''\to X'\to X$，其中
$X'\to X$ 是 $X$ 在 $x$ 处的爆破
（定义 \ref{spaces-resolve-definition-blowup-at-point}），
而 $X''\to X'$ 是 $X'$ 的正规化。
\end{definition}

\noindent
这里，正规化 $X''\to X'$ 作为代数空间有定义，% P10-E0410
因为 $X'$ 满足《空间的态射》引理
\ref{spaces-morphisms-lemma-prepare-normalization} 中的等价条件；
这由《空间上的除子》引理
\ref{spaces-divisors-lemma-blowup-finite-nr-irreducibles} 得到。
正规化的定义见《空间的态射》定义
\ref{spaces-morphisms-definition-normalization}。

\medskip\noindent
一般而言，即使 $X$ 诺特，正规化爆破也未必固有。
回忆，若一个代数空间有由 Nagata 环的谱所成的仿射平展覆盖，
则称它为 Nagata；见《空间的性质》定义
\ref{spaces-properties-definition-type-property}、
备注 \ref{spaces-properties-remark-list-properties-local-etale-topology}
以及《性质》定义 \ref{properties-definition-nagata}。

\begin{lemma}
\label{spaces-resolve-lemma-Nagata-normalized-blowup}
在定义 \ref{spaces-resolve-definition-normalized-blowup} 中，若 $X$ 为 Nagata，
则 $X$ 在 $x$ 处的正规化爆破是 $X$ 上固有的正规 Nagata 代数空间。
\end{lemma}

\begin{proof}
爆破态射 $X'\to X$ 固有（由于 $X$ 局部诺特，可以应用
《空间上的除子》引理
\ref{spaces-divisors-lemma-blowing-up-projective}）。
因此 $X'$ 为 Nagata（《空间的态射》引理
\ref{spaces-morphisms-lemma-finite-type-nagata}）。
故正规化 $X''\to X'$ 有限（《空间的态射》引理
\ref{spaces-morphisms-lemma-nagata-normalization}），
从而 $X''\to X$ 也固有（《空间的态射》引理
\ref{spaces-morphisms-lemma-finite-proper} 和
\ref{spaces-morphisms-lemma-composition-proper}）。
由此，正规化爆破是正规（《空间的态射》引理
\ref{spaces-morphisms-lemma-normalization-normal}）Nagata 代数空间。
\end{proof}

\noindent
下面是引理 \ref{spaces-resolve-lemma-equivalence-sequence-blowups}
对正规化爆破的类似结论。

\begin{lemma}
\label{spaces-resolve-lemma-equivalence-sequence-normalized-blowups}
设 $X,x_i,U_i\to X,u_i$ 如 (\ref{spaces-resolve-equation-equivalence}) 中所述，
并假设在 $F$ 下 $f:Y\to X$ 对应于 $g_i:Y_i\to U_i$。
假设 $X$ 满足《空间的态射》引理
\ref{spaces-morphisms-lemma-prepare-normalization} 中的等价条件。
则 $f$ 存在分解
$$
Y = Z_m \to Z_{m - 1} \to \ldots \to Z_1 \to Z_0 = X
$$
其中 $Z_{j+1}\to Z_j$ 是 $Z_j$ 在位于
$\{x_1,\ldots,x_n\}$ 之上的闭点 $z_j$ 处的正规化爆破，
当且仅当对每个 $i$，$g_i$ 存在分解
$$
Y_i = Z_{i,m_i} \to Z_{i,m_i-1} \to \ldots \to Z_{i,1}
\to Z_{i,0}=U_i
$$
其中 $Z_{i,j+1}\to Z_{i,j}$ 是 $Z_{i,j}$
在位于 $u_i$ 之上的闭点 $z_{i,j}$ 处的正规化爆破。
\end{lemma}

\begin{proof}
这由证明引理 \ref{spaces-resolve-lemma-equivalence-sequence-blowups}
时完全相同的论证得出。
\end{proof}

\noindent
Nagata 代数空间局部诺特。

\begin{lemma}
\label{spaces-resolve-lemma-dominate-by-normalized-blowing-up}
设 $S$ 为概形，$X$ 为 $S$ 上维数 $\dim(X)=2$ 的诺特
Nagata 代数空间，$f:Y\to X$ 为固有双有理态射。
则存在交换图
$$
\xymatrix{
X_n \ar[r] \ar[d] &
X_{n - 1} \ar[r] &
\ldots \ar[r] &
X_1 \ar[r] &
X_0 \ar[d] \\
Y \ar[rrrr]  & & & & X
}
$$
其中 $X_0\to X$ 是正规化，而
$X_{i+1}\to X_i$ 是 $X_i$ 在某个闭点处的正规化爆破。
\end{lemma}

\begin{proof}
虽然可以直接为代数空间证明本引理，
我们仍沿用上面的做法，将它约化到概形情形。

\medskip\noindent
我们将用到：诺特代数空间拟分离，因此其点有良定义的剩余域
（例如由《合宜空间》引理
\ref{decent-spaces-lemma-decent-space-elementary-etale-neighbourhood}）。
我们还将不再另行说明地使用《空间的态射》第
\ref{spaces-morphisms-section-nagata}、
\ref{spaces-morphisms-section-dimension-formula} 和
\ref{spaces-morphisms-section-normalization} 节的结果。
可用 $Y$ 的正规化替换 $Y$。令 $X_0\to X$ 为正规化。
态射 $Y\to X$ 分解经过 $X_0$，故可以假设 $X$ 与 $Y$ 都正规。

\medskip\noindent
假设 $X$ 与 $Y$ 正规。态射 $f:Y\to X$ 在一个开集上为同构，
该开集包含 $Y$ 中每个余维 $0$ 或 $1$ 的点，以及纤维为有限的每个点；
见《域上的空间》引理
\ref{spaces-over-fields-lemma-modification-normal-iso-over-codimension-1}。
由此存在闭点的有限集合 $T\subset|X|$，使 $f$
在 $X\setminus T$ 上为同构。
由《空间态射进阶》引理
\ref{spaces-more-morphisms-lemma-dominate-modification-by-blowup}，
存在支配 $Y$ 的 $X\setminus T$-容许爆破 $Y'\to X$。
用 $Y'$ 的正规化替换 $Y$ 后，可以假设 $Y\to X$ 可表。

\medskip\noindent
设 $T=\{x_1,\ldots,x_r\}$。如第 \ref{spaces-resolve-section-strategy} 节，
选取初等平展邻域 $(U_i,u_i)\to(X,x_i)$。
对每个 $i$，态射 $Y_i=Y\times_XU_i\to U_i$ 是固有双有理态射，
并且在 $U_i\setminus\{u_i\}$ 上为同构。
因此可以应用《曲面解消》引理
\ref{resolve-lemma-dominate-by-normalized-blowing-up}，得到序列
$$
X_{i,m_i}\to X_{i,m_i-1}\to\ldots\to X_1\to X_{i,0}=U_i
$$
% P10-E0411
其中各态射是在位于 $u_i$ 之上的闭点处作正规化爆破，
且 $X_{i,m_i}$ 支配 $Y_i$。
由引理 \ref{spaces-resolve-lemma-equivalence-sequence-normalized-blowups}，
得到一列正规化爆破
$$
X_m\to X_{m-1}\to\ldots\to X_1\to X_0=X
$$
如引理所述，其到各 $U_i$ 的基变换给出上述序列。
由引理 \ref{spaces-resolve-lemma-equivalence} 的范畴等价，$X_m$ 支配 $Y$。
\end{proof}















\section{到完备化的基变换}
\label{spaces-resolve-section-aux}

\noindent
下面这个简单引理将在后文中成为有用的工具。

\begin{lemma}
\label{spaces-resolve-lemma-iso-completions}
设 $(A,\mathfrak m,\kappa)$ 为局部环，极大理想 $\mathfrak m$
有限生成。设 $X$ 为 $A$ 上的合宜代数空间，并令
$Y=X\times_{\Spec(A)}\Spec(A^\wedge)$，其中 $A^\wedge$
为 $A$ 的 $\mathfrak m$-进完备化。
设 $q\in|Y|$ 的像 $p\in|X|$ 位于 $\Spec(A)$ 的闭点之上，
则 Hensel 局部环的映射
$\mathcal{O}_{X,p}^h\to\mathcal{O}_{Y,q}^h$
在完备化上诱导同构。
\end{lemma}

\begin{proof}
选取《合宜空间》引理
\ref{decent-spaces-lemma-decent-space-elementary-etale-neighbourhood}
中的初等平展邻域 $(U,u)\to(X,p)$，并令
$V=U\times_XY=U\times_{\Spec(A)}\Spec(A^\wedge)$、$v=(u,q)$，
便立刻从概形情形得到结论。概形情形见《曲面解消》引理
\ref{resolve-lemma-iso-completions}。
\end{proof}

\begin{lemma}
\label{spaces-resolve-lemma-port-regularity-to-completion}
设 $(A,\mathfrak m,\kappa)$ 为诺特局部环。
设 $X\to\Spec(A)$ 为局部有限型态射，且 $X$ 是合宜代数空间。
令 $Y=X\times_{\Spec(A)}\Spec(A^\wedge)$。
设 $y\in|Y|$ 的像为 $x\in|X|$。则：
\begin{enumerate}
\item 若 $\mathcal{O}_{Y,y}^h$ 正则，则 $\mathcal{O}_{X,x}^h$ 正则；
\item 若 $y$ 在闭纤维中，则 $\mathcal{O}_{Y,y}^h$ 正则
$\Leftrightarrow\mathcal{O}_{X,x}^h$ 正则；
\item 若 $X$ 在 $A$ 上固有，则 $X$ 正则当且仅当 $Y$ 正则。
\end{enumerate}
\end{lemma}

\begin{proof}
经过平展局部化，前两个断言立即归结为相应的概形结论；见《曲面解消》引理
\ref{resolve-lemma-port-regularity-to-completion}。
对 (3)，由于 $Y\to X$ 满（因为 $A\to A^\wedge$ 忠实平坦），
由 (1) 可知 $Y$ 正则蕴含 $X$ 正则。
反之，若 $X$ 正则，则 $Y$ 在特殊纤维所有点处的 Hensel 局部环都正则。
任取 $y\in|Y|$。在固有情形中，
$|Y|\to|\Spec(A^\wedge)|$ 是闭映射，故可找到特化
$y\leadsto y_0$，其中 $y_0$ 在闭纤维中。
依照《合宜空间》引理
\ref{decent-spaces-lemma-decent-space-elementary-etale-neighbourhood}，
选取初等平展邻域 $(V,v_0)\to(Y,y_0)$。
由于 $Y$ 合宜，可以把 $y\leadsto y_0$ 提升为 $V$ 中的特化
$v\leadsto v_0$（《合宜空间》引理
\ref{decent-spaces-lemma-decent-specialization}）。
于是 $\mathcal{O}_{V,v}$ 是 $\mathcal{O}_{V,v_0}$ 的局部化，故正则，
证明完成。
\end{proof}

\begin{lemma}
\label{spaces-resolve-lemma-formally-unramified}
设 $(A,\mathfrak m)$ 为局部诺特环，$X$ 为 $A$ 上的代数空间。
假设：
\begin{enumerate}
\item $A$ 解析无分歧
（《代数》定义 \ref{algebra-definition-analytically-unramified}）；
\item $X$ 在 $A$ 上局部有限型；
\item $X\to\Spec(A)$ 在 $X$ 中每个余维 $0$ 的点处平展。
\end{enumerate}
则 $X$ 的正规化在 $X$ 上有限。
\end{lemma}

\begin{proof}
选取概形 $U$ 及满平展态射 $U\to X$。
于是 $U\to\Spec(A)$ 满足《曲面解消》引理
\ref{resolve-lemma-formally-unramified} 的假设，因而也满足其结论。
\end{proof}











\section{推得的性质}
\label{spaces-resolve-section-existence-gives}

\noindent
本节证明：对诺特整代数空间，正则变换的存在会带来相当多的后果。
对“病态”诺特代数空间不感兴趣的读者可以跳过本节。

\begin{lemma}
\label{spaces-resolve-lemma-regular-alteration-implies}
设 $S$ 为概形，$Y$ 为 $S$ 上的诺特整代数空间。
假设存在变换 $f:X\to Y$，且 $X$ 正则。
则正规化 $Y^\nu\to Y$ 有限，并且 $Y$ 有一个正则的稠密开子空间。
\end{lemma}

\begin{proof}
通过平展局部化，只须在 $Y=\Spec(A)$ 时证明，其中 $A$ 是诺特整环。
设 $B$ 为 $A$ 在其分式域中的整闭包，并令
$C=\Gamma(X,\mathcal{O}_X)$。由《空间的上同调》引理
\ref{spaces-cohomology-lemma-proper-pushforward-coherent}，
$C$ 是有限 $A$-模。
由于 $X$ 正规（《空间的性质》引理
\ref{spaces-properties-lemma-regular-normal}），
$C$ 是正规整环（《域上的空间》引理
\ref{spaces-over-fields-lemma-normal-integral-sections}）。
因此 $B\subset C$；又因 $A$ 诺特，得出 $B$ 在 $A$ 上有限。

\medskip\noindent
存在非空开集 $V\subset Y$，使 $f^{-1}V\to V$ 有限；
见《域上的空间》定义 \ref{spaces-over-fields-definition-alteration}。
缩小 $V$ 后，可以假设 $f^{-1}V\to V$ 平坦
（《空间的态射》命题
\ref{spaces-morphisms-proposition-generic-flatness-reduced}）。
因此 $f^{-1}V\to V$ 忠实平坦。
由《代数》引理 \ref{algebra-lemma-descent-regular}，$V$ 正则。
\end{proof}

\begin{lemma}
\label{spaces-resolve-lemma-regular-alteration-implies-local}
设 $(A,\mathfrak m,\kappa)$ 为局部诺特整环。
假设存在变换 $f:X\to\Spec(A)$，且 $X$ 正则。则：
\begin{enumerate}
\item 存在非零 $f\in A$，使 $A_f$ 正则；
\item $A$ 在其分式域中的整闭包 $B$ 在 $A$ 上有限；
\item $B$ 的 $\mathfrak m$-进完备化是正规环；换言之，
$B$ 在各极大理想处的完备化都是正规整环；
\item $A$ 的泛形式纤维正则。
\end{enumerate}
\end{lemma}

\begin{proof}
(1) 与 (2) 由引理 \ref{spaces-resolve-lemma-regular-alteration-implies} 得出。
为给 (3) 的证明建立记号，需要重做该引理证明的一部分。
令 $C=\Gamma(X,\mathcal{O}_X)$。由《空间的上同调》引理
\ref{spaces-cohomology-lemma-proper-pushforward-coherent}，
$C$ 是有限 $A$-模。
由于 $X$ 正规（《空间的性质》引理
\ref{spaces-properties-lemma-regular-normal}），
$C$ 是正规整环（《域上的空间》引理
\ref{spaces-over-fields-lemma-normal-integral-sections}）。
因此 $B\subset C$；又因 $A$ 诺特，$B$ 在 $A$ 上有限。
由《曲面解消》引理 \ref{resolve-lemma-algebra-helper}，
为证明 (3)，只须证明 $C$ 的 $\mathfrak m$-进完备化 $C^\wedge$ 正规。

\medskip\noindent
由《代数》引理 \ref{algebra-lemma-completion-finite-extension}，
完备化 $C^\wedge$ 是 $C$ 在位于 $\mathfrak m$ 之上的各素理想处
完备化的乘积。这些素理想为有限多个，且就是 $C$ 的极大理想
$\mathfrak m_1,\ldots,\mathfrak m_r$。
（对 $B$ 的相应结果解释了引理中的最后一个说明。）
因此，用 $C_{\mathfrak m_i}$ 替换 $A$，并用
$X_i=X\times_{\Spec(C)}\Spec(C_{\mathfrak m_i})$ 替换 $X$，
便约化到下一段讨论的情形。
（注意，由《空间的上同调》引理
\ref{spaces-cohomology-lemma-flat-base-change-cohomology}，
$\Gamma(X_i,\mathcal{O})=C_{\mathfrak m_i}$。）

\medskip\noindent
现在 $A$ 是诺特局部正规整环，$f:X\to\Spec(A)$
是正则变换，并且 $\Gamma(X,\mathcal{O}_X)=A$。
我们须证明 $A$ 的完备化 $A^\wedge$ 是正规整环。
由引理 \ref{spaces-resolve-lemma-port-regularity-to-completion}，
$Y=X\times_{\Spec(A)}\Spec(A^\wedge)$ 正则。
由《空间的上同调》引理
\ref{spaces-cohomology-lemma-flat-base-change-cohomology}，
$\Gamma(Y,\mathcal{O}_Y)=A^\wedge$。% P10-E0412
与前面相同，可推出 $A^\wedge$ 正规。
具体而言，由《空间的性质》引理
\ref{spaces-properties-lemma-regular-normal}，$Y$ 正规。
它连通，因为 $\Gamma(Y,\mathcal{O}_Y)=A^\wedge$ 是局部环。
所以 $Y$ 正规且整（对分离代数空间，连通且正规蕴含整）。
于是由《域上的空间》引理
\ref{spaces-over-fields-lemma-normal-integral-sections}，
$\Gamma(Y,\mathcal{O}_Y)=A^\wedge$ 是正规整环。这证明了 (3)。

\medskip\noindent
证明 (4)。令 $\eta\in\Spec(A)$ 表示泛点，并以指标 $\eta$
表示到 $\eta$ 的基变换。由于 $f$ 是变换，概形 $X_\eta$
在 $\eta$ 上有限且忠实平坦。由引理
\ref{spaces-resolve-lemma-port-regularity-to-completion}，
$Y=X\times_{\Spec(A)}\Spec(A^\wedge)$ 正则，
故 $Y_\eta$ 正则（它是 $Y$ 中开集的极限）。
于是
$Y_\eta\to\Spec(A^\wedge\otimes_A\kappa(\eta))$
有限且忠实平坦地映到泛形式纤维。
结论由《代数》引理 \ref{algebra-lemma-descent-regular} 得出。
\end{proof}












\section{解消}
\label{spaces-resolve-section-resolution}

\noindent
先给出一个定义。

\begin{definition}
\label{spaces-resolve-definition-resolution}
设 $S$ 为概形，$Y$ 为 $S$ 上的诺特整代数空间。
所谓 $X$ 的一个{\it 奇点解消}，% P10-E0413
是修改 $f:X\to Y$，其中 $X$ 正则。
\end{definition}

\noindent
对曲面，我们有时还需要多一点信息。

\begin{definition}
\label{spaces-resolve-definition-resolution-surface}
设 $S$ 为概形，$Y$ 为 $S$ 上维数为 $2$ 的诺特整代数空间。
若存在序列
$$
Y_n \to X_{n - 1} \to \ldots \to Y_1 \to Y_0 \to Y
$$
% P10-E0414
满足
\begin{enumerate}
\item 对 $i=0,\ldots,n$，$Y_i$ 在 $Y$ 上固有；
\item $Y_0\to Y$ 是正规化；
\item 对 $i=1,\ldots,n$，$Y_i\to Y_{i-1}$ 是正规化爆破；
\item $Y_n$ 正则，
\end{enumerate}
则称 $Y$ 有一个{\it 由正规化爆破给出的奇点解消}。
\end{definition}

\noindent
注意，条件 (1) 蕴含 $Y$ 的正规化 $Y_0$ 在 $Y$ 上有限，
并且正规化爆破中所用的各正规化也都是有限的。
现在终于来到本章的主定理。

\begin{theorem}
\label{spaces-resolve-theorem-resolve}
设 $S$ 为概形，$Y$ 为 $S$ 上的二维整诺特代数空间。
下列条件等价：
\begin{enumerate}
\item 存在变换 $X\to Y$，其中 $X$ 正则；
\item $Y$ 存在奇点解消；
\item $Y$ 有由正规化爆破给出的奇点解消；
\item 正规化 $Y^\nu\to Y$ 有限，$Y^\nu$ 有有限多个奇点
$y_1,\ldots,y_m\in|Y|$，% P10-E0415
并且对每个 $i$，Hensel 局部环
$\mathcal{O}_{Y^\nu,y_i}^h$ 的完备化正规。
\end{enumerate}
\end{theorem}

\begin{proof}
蕴含关系 (3) $\Rightarrow$ (2) $\Rightarrow$ (1) 是立即的。

\medskip\noindent
设 $X\to Y$ 为变换且 $X$ 正则。
由引理 \ref{spaces-resolve-lemma-regular-alteration-implies}，$Y^\nu\to Y$ 有限。
考虑《空间的态射》引理
\ref{spaces-morphisms-lemma-normalization-normal}
给出的分解 $f:X\to Y^\nu$。
由《域上的空间》引理
\ref{spaces-over-fields-lemma-finite-in-codim-1}，
$f$ 在开集 $V\subset Y^\nu$ 上有限，而该开集包含 $Y^\nu$
中每个余维 $\leq1$ 的点。
由《代数》引理 \ref{algebra-lemma-CM-over-regular-flat}，
以及维数 $\leq 2$ 的正规局部环由 Serre 判别准则成为 Cohen--Macaulay 环这一事实
（《代数》引理 \ref{algebra-lemma-criterion-normal}），
$f$ 在 $V$ 上平坦。
再由《代数》引理 \ref{algebra-lemma-descent-regular}，$V$ 正则。
由于 $Y^\nu$ 诺特，得出
$Y^\nu\setminus V=\{y_1,\ldots,y_m\}$ 有限。
对每个 $i$，令 $\mathcal{O}_{Y^\nu,y_i}^h$ 为 Hensel 局部环。
则
$X\times_Y\Spec(\mathcal{O}_{Y^\nu,y_i}^h)$
是 $\Spec(\mathcal{O}_{Y^\nu,y_i}^h)$ 的正则变换
（省略一些细节）。由引理
\ref{spaces-resolve-lemma-regular-alteration-implies-local}，
$\mathcal{O}_{Y^\nu,y_i}^h$ 的完备化正规。
这就证明了 (1) $\Rightarrow$ (4)。

\medskip\noindent
假设 (4)，须证明 (3)。可以立刻用其正规化替换 $Y$。
设 $y_1,\ldots,y_m\in|Y|$ 为奇点。
如第 \ref{spaces-resolve-section-strategy} 节，选取一族初等平展邻域
$(V_i,v_i)\to(Y,y_i)$。
对每个 $i$，Hensel 局部环 $\mathcal{O}_{Y^\nu,y_i}^h$
是 $\mathcal{O}_{V_i,v_i}$ 的 Hensel 化，故这些环的完备化同构。
于是由概形情形的结果（《曲面解消》定理
\ref{resolve-theorem-resolve}），存在有限列正规化爆破
$$
X_{i,n_i}\to X_{i,n_i-1}\to\ldots\to V_i
$$
其中只爆破位于 $v_i$ 之上的点，并且 $X_{i,n_i}$ 正则。
由引理 \ref{spaces-resolve-lemma-equivalence-sequence-normalized-blowups}，
存在一列正规化爆破
$$
X_n\to X_{n-1}\to\ldots\to X_1\to Y
$$
当然 $X_n$ 也正则（考察局部环即可）。证明完成。
\end{proof}















\section{例子}
\label{spaces-resolve-section-examples}

\noindent
下面给出一些与本章前述结果有关的例子。

\begin{example}
\label{spaces-resolve-example-factorial}
\begin{reference}
\cite[4(c)]{Samuel-UFD}
\end{reference}
设 $k$ 为域。若 $r,s,t$ 两两互素，则环
$A=k[x,y,z]/(x^r+y^s+z^t)$ 是唯一分解整环（UFD）。
事实上，由于 $x^r+y^s+z^t$ 不可约，$A$ 是整环。
元素 $z$ 是素元，即它在 $A$ 中生成素理想。
另一方面，若对某个 $e$ 有 $t=1+ers$，则
$$
A[1/z]\cong k[x',y',1/z]
$$
其中 $x'=x/z^{es}$、$y'=y/z^{er}$，并且
$z=(x')^r+(y')^s$。
因此 $A[1/z]$ 是多项式环的一个局部化，故为 UFD。
由 Nagata 的一个论证，$A$ 也是 UFD；见《代数》引理
\ref{algebra-lemma-invert-prime-elements}。
当 $t$ 模 $rs$ 不同余于 $1$ 时，也可以给出类似论证。
\end{example}

\begin{example}
\label{spaces-resolve-example-completion-not-factorial}
\begin{reference}
一般情形中局部 Picard 群的非零性见
\cite{Brieskorn} 和 \cite{Lipman-rational}。
\end{reference}
若 $1<r<s<t$ 是两两互素且不等于 $2,3,5$ 的整数，则环
$A=\mathbf{C}[[x,y,z]]/(x^r+y^s+z^t)$ 不是 UFD。
例如考虑特殊情形
$A=\mathbf{C}[[x,y,z]]/(x^2+y^5+z^7)$。
考虑映射
$$
\psi_\zeta:\mathbf{C}[[x,y,z]]/(x^2+y^5+z^7)\to\mathbf{C}[[t]]
$$
其定义为
$$
x\mapsto t^7,\quad
y\mapsto t^3,\quad
z\mapsto-\zeta t^2(1+t)^{1/7}
$$
其中 $\zeta$ 是 $7$ 次单位根。
$\psi_\zeta$ 的核 $\mathfrak p_\zeta$ 是高度一素理想；
因而若 $A$ 是 UFD，它就是主理想，不妨设由
$f_\zeta\in\mathbf{C}[[x,y,z]]$ 生成。
注意 $V(x^3-y^7)=\bigcup V(\mathfrak p_\zeta)$，并且
$A/(x^3-y^7)$ 在闭点以外是既约的。
因此，仍假设 $A$ 为 UFD，就会有
$$
\prod\nolimits_\zeta f_\zeta
=u(x^3-y^7)+a(x^2+y^5+z^7)
\quad\text{于}\quad
\mathbf{C}[[x,y,z]]
$$
其中 $u\in\mathbf{C}[[x,y,z]]$ 是单位，
$a\in\mathbf{C}[[x,y,z]]$。
乘以一个常数后，可以假设 $u(0,0,0)=1$。
注意左端的消失阶为 $7$，故 $a=-x\bmod\mathfrak m^2$。
但这样一来，右端出现项 $xy^5$，而左端没有该项，矛盾。
\end{example}

\begin{example}
\label{spaces-resolve-example-not-blow-up}
存在一个优秀的维数为 $2$ 的诺特局部环，以及一个不是概形的修改
$X\to S=\Spec(A)$。我们概述一种构造。
设 $X$ 为 $\mathbf{C}$ 上的正规曲面，只有一个奇点 $x\in X$。
假设存在解消 $\pi:X'\to X$，使例外纤维
$C=\pi^{-1}(x)_{red}$ 是光滑射影曲线。
再假设存在 $c\in C$，使得只要 $\mathcal{O}_C(nc)$
在 $\Pic(X')\to\Pic(C)$ 的像中，就有 $n=0$。
令 $X''\to X'$ 为在非奇异点 $c$ 处的爆破。
设 $C'\subset X''$ 为 $C$ 的严格变换，$E\subset X''$ 为例外纤维。
由 Artin 的结果（\cite{ArtinII}；例如用 \cite{Mumford-topology}
可见 $C'$ 的法丛为负），可以在 $X''$ 中收缩曲线 $C'$，
得到代数空间 $X'''$。图示为
$$
\xymatrix{
& X'' \ar[ld] \ar[rd] \\
X' \ar[rd] &  & X''' \ar[ld] \\
& X
}
$$
我们断言 $X'''$ 不是概形。这给出了所需例子，
因为 $X'''$ 是概形，当且仅当 $X'''$ 到
$A=\mathcal{O}_{X,x}$ 的基变换是概形（省略细节）。
如果 $X'''$ 是概形，% P10-E0416
则 $C'$ 在 $X'''$ 中的像会有一个仿射邻域。
该邻域的补集会是 $X'''$ 上的有效 Cartier 除子
（因为 $X'''$ 除 $1$ 个点外非奇异）。
这个有效 Cartier 除子对应于 $X''$ 上一个与 $E$ 相交且避开 $C'$
的有效 Cartier 除子。取其在 $X'$ 中的像，就得到一个有效
Cartier 除子，它在集合论意义下与 $C$ 交于 $c$。
这不可能，因为按假设，$c$ 的任何倍数都不是某个 Cartier 除子的限制。

\medskip\noindent
最后还须找到这样的奇异曲面 $X$。可以直接取 $X$ 为
$\mathbf{A}^3_\mathbf{C}=\Spec(\mathbf{C}[x,y,z])$ 中由
$$
x^3+y^3+z^3+x^4+y^4+z^4=0
$$
给出的仿射曲面，奇点为 $(0,0,0)$。
此时 $(0,0,0)$ 是唯一奇点。
沿对应于 $(0,0,0)$ 的极大理想爆破 $X$，得到三个图册，
每个都同构于光滑仿射曲面
$$
1+s^3+t^3+x(1+s^4+t^4)=0
$$
它非奇异，例外除子 $C$ 由 $x=0$ 给出。
读者会认出 $C$ 是椭圆曲线。
此外，从 $(0,0,0)$ 作投影可见曲面 $X$ 是有理的；
也可在爆破方程中解出 $x$ 得到同一结论。
最后，非奇异有理曲面的 Picard 群可数，
而复数域上椭圆曲线的 Picard 群不可数。
故可以找到所要求的闭点 $c$。
\end{example}
